\scanpage{71}
% archive sheet number: 70
% @ L1–9, crossed continuation at top
\deleted{cas au moins, c'est dû au fait que je n'ai pas tenu compte des arcs de spécialisations exceptionnels (vers leurs pos. supersingulières). Là encore, il y a deux variantes, \unsure{suivant|suivt} qu'on se permet de se spécialiser en une pos. rel supersingulière, ou qu'on veut l'éviter, en passant \deleted{directement} \unsure{à} la position voisine opposée}

% @ L10–18
Je viens de vérifier, sauf erreur (j'ai fait pas mal d'erreurs encore dans le cours de la vérification !) les trois procédés $A_1,A_2,B_{\unsure{1|3}}$ donnent lieu aux bons circuits \margin{\unsure{--\,.}}autour d'une position supersingulière, il n'y a que $B_2$ qui déconne complètement (or il semblerait que c'était sur $B_2$ que je m'étais arrêté pendant un jour ou deux, Dieu sait pourquoi !)

% @ L19–22
Le plus intéressant, après $A_2$, semble bien sûr $A_1$, mais je viens encore d'essayer de comprendre les circuits auxquels il donne naissance, sans y réussir.
